\chapter{Conclusions}
\label{chap:conclusions}

This thesis has presented the design, implementation, and evaluation of a
complete Retrieval-Augmented Generation system built on top of a vector
database. The system allows authenticated users to upload PDF documents,
ask questions in natural language, and receive grounded answers together
with explicit source citations down to the page level. This final chapter
summarises the contributions, discusses the limitations that were revealed
by the evaluation, and outlines several avenues for future work.

\section{Summary of Contributions}
\label{sec:summary-contributions}

The main contributions of the thesis are the following.

\textbf{A focused, layered reference implementation.} The backend is
organised into thin API routes, services, repositories, and core utilities,
with strict boundaries enforced by dependency injection. The structure
makes every responsibility visible and the system easy to reason about ---
a deliberate choice for a work that is meant to double as a teaching
example. The entire backend is test-driven, with 22 tests passing and
83\,\% branch coverage; the authentication layer in particular is at
96\,\% coverage, with end-to-end integration tests for register, login,
and protected routes.

\textbf{An end-to-end RAG pipeline.} The ingestion pipeline handles PDF
uploads with PyMuPDF, produces overlapping chunks of approximately 500
tokens, computes dense embeddings with Google Gemini, and stores them in
a per-user ChromaDB collection along with the metadata needed for
citation (filename, page, document id). The query pipeline embeds the
question, retrieves the top-5 chunks, builds a grounded prompt, calls
Gemini, and returns the answer together with machine-readable citations.
Every stage of the pipeline is mockable for unit testing.

\textbf{Strict per-user isolation.} Isolation rests on two independent
mechanisms: JWT-based authentication at the HTTP layer, and per-user
collections at the vector-store layer. The second mechanism is
architectural rather than filter-based --- a user cannot accidentally
read another user's chunks even if a bug were to bypass the HTTP-level
check --- which makes the safe behaviour the default.

\textbf{A production-ready deployment.} The entire stack is packaged as
a Docker Compose project with three services (PostgreSQL, FastAPI
backend, and nginx serving the React frontend). A single \texttt{docker
compose up --build} command brings up the complete system. The
PostgreSQL container is exposed on a non-standard host port (5433) to
avoid conflicts with native installs, and a single \texttt{.env} file
drives all configuration.

\textbf{A modern frontend.} The React + Vite + TypeScript + Tailwind +
shadcn/ui frontend provides dark-mode support out of the box, a
themeable component library, and a dedicated help page written in
Romanian for end users. Source citations are rendered as clickable
badges beneath every answer.

\textbf{An experimental evaluation.} The evaluation in
Chapter~\ref{chap:evaluation} measured retrieval quality, answer
groundedness, and latency on two representative corpora, and compared
the dense-retrieval component against a PostgreSQL full-text search
baseline on exactly the same data. Dense retrieval reached 0.93 and
0.97 recall@5 on the legal and restaurant corpora respectively,
outperforming the lexical baseline by 0.16 and 0.04 points. Grounded
answer rates were 0.90 and 0.97, and the median query latency was
1.6\,s, almost entirely consumed by the LLM call.

\section{Limitations}
\label{sec:limitations}

The evaluation and the implementation also exposed several limitations
that a production deployment would need to address.

\textbf{Per-user collections scale linearly.} Using one ChromaDB
collection per user is simple and safe but becomes expensive at very
large user counts: every collection carries a small fixed overhead in
ChromaDB, and the directory on disk grows with the number of
collections. A multi-tenant deployment with thousands of users would
likely need a single collection partitioned by user id as a metadata
filter, trading some of the safety guarantees of the current design
for better scaling.

\textbf{Free-tier rate limits dominate latency.} The query latency
reported in Chapter~\ref{chap:evaluation} is almost entirely spent in
the Gemini generation call, and the free-tier quota constrains both
throughput and tail latencies. Moving to a paid tier, or hosting a
smaller open-weights model locally, would be a necessary step for
production traffic.

\textbf{Retrieval relies on a single encoder.} The system uses the
Gemini \texttt{text-embedding-004} model for both indexing and query
embedding. The encoder was trained primarily on English and
well-represented multilingual corpora; for more specialised domains, a
domain-adapted or fine-tuned encoder would likely improve retrieval
further, but fine-tuning against Gemini is not available and would
require migrating to an open encoder (for example a Sentence-BERT
variant \cite{Reimers2019SentenceBERT}).

\textbf{No document metadata in PostgreSQL.} Documents live entirely in
ChromaDB, which is sufficient for the present scope but prevents
features such as server-side search over filenames, document sharing
between users, or audit logs. Adding a \texttt{documents} table in
PostgreSQL would be a small change but was deferred to keep the schema
minimal.

\textbf{No OCR.} PDFs with scanned pages but no text layer are silently
skipped. A practical deployment would need to detect them and run them
through an OCR stage (for example Tesseract or a cloud OCR API) before
ingestion.

\textbf{Synchronous ingestion.} PDF uploads are processed synchronously
inside the HTTP request. For large documents this can exceed typical
HTTP timeouts. Offloading ingestion to a background worker queue is a
natural next step.

\section{Future Work}
\label{sec:future-work}

The architecture leaves room for several improvements that would be
good candidates for follow-up projects:

\textbf{Hybrid retrieval.} As the evaluation showed, a small number of
failures at the retrieval step are lexically obvious but semantically
distant. A reranking stage that combines the dense top-$k$ with a BM25
score would likely close most of this gap.

\textbf{Query rewriting.} When a user asks a long or ambiguous question,
retrieval quality can be improved by first rewriting the question into
a more retrieval-friendly form using the LLM, and then embedding the
rewritten form. This is cheap in comparison to the main generation
call and has been reported to help in several open-source RAG
frameworks.

\textbf{Streaming responses.} The Gemini API supports token-level
streaming, and FastAPI can forward a streaming response to the client.
Implementing this would substantially improve the user-perceived
latency, since the first token typically arrives in under 500\,ms.

\textbf{Multi-tenant deployment.} Moving from per-user collections to a
single collection with metadata-based filtering, as discussed above,
would be necessary for scaling beyond a few hundred users. It would
require a rigorous review of the authorisation paths to preserve the
isolation guarantees that the current architecture enforces structurally.

\textbf{Wider evaluation.} The present evaluation is limited to two
corpora and 60 hand-written questions. A larger-scale evaluation on
public benchmarks such as NaturalQuestions or HotpotQA, combined with
automated grounding metrics, would give more robust numbers, although
such benchmarks are in English and do not directly reflect the use
cases targeted by the system.

\vspace{1em}

In conclusion, the thesis demonstrates that a modest but carefully
engineered stack --- a layered FastAPI backend, a simple embedded
vector database, a pre-trained commercial LLM, and a typed React
frontend --- is sufficient to build a grounded, per-user question
answering system that outperforms lexical search on realistic
corpora. The same stack can serve as a starting point for further
research and for real-world deployments, and it remains small enough
to be understood in its entirety by a single reader.
